%%%%%%%%%%%%%%%%%%%%%%%%%%%%%%%%%%%%%%%%%%%%%%%%%%%%%%%%%%%%%%%%%%%%%%%%%%%%
% Supporting Information for the WRR manuscript
% Compile main.tex first so that xr can import its labels.
%%%%%%%%%%%%%%%%%%%%%%%%%%%%%%%%%%%%%%%%%%%%%%%%%%%%%%%%%%%%%%%%%%%%%%%%%%%%
\documentclass[draft,wrr]{agutexSI2019}

\ExplSyntaxOn
\msg_redirect_name:nnn { mark } { infinite-shrinkage } { none }
\ExplSyntaxOff

\usepackage{amsmath,amssymb}
\usepackage{array}
\usepackage{booktabs}
\usepackage{graphicx}
\usepackage{lineno}
\usepackage{xcolor}
\usepackage{xr}
\setkeys{Gin}{draft=false}
\externaldocument{main}

\renewcommand{\theequation}{S\arabic{equation}}
\newcounter{sitext}
\renewcommand{\thesitext}{S\arabic{sitext}}
\newcommand{\sitext}[1]{%
  \refstepcounter{sitext}%
  \par\medskip\noindent\textbf{Text \thesitext. #1}\par\medskip}
\newcommand{\sisub}[1]{\par\medskip\noindent\textbf{#1}\par\nobreak\smallskip}
\newcommand{\verify}[1]{%
  \par\smallskip\noindent\textcolor{red}{\textbf{Author query---}#1}\par\smallskip}
\newcommand{\NaIso}{\(^{23}\mathrm{Na}\)}
\newcommand{\HIso}{\(^{1}\mathrm{H}\)}
\linenumbers

\begin{document}
\title{Supporting Information for: Evaporation-Driven Solute Redistribution
in a Macroscopically Heterogeneous Porous Medium With a Vertical Textural
Interface: MRI Observations and Numerical Analysis}

\makeatletter
\def\@oddfoot{}
\def\@evenfoot{}
\makeatother

\authors{M. A. Chaudhry\affil{2}, Z. Chen\affil{2}\thanks{Corresponding author:
zhixin.chen@iws.uni-stuttgart.de}, A. Pohlmeier\affil{1}, R. Helmig\affil{2},
and J. A. Huisman\affil{1}}

\affiliation{1}{Agrosphere Institute (IBG 3), Forschungszentrum J\"ulich GmbH,
52428 J\"ulich, Germany.}
\affiliation{2}{Institute for Modeling Hydraulic and Environmental Systems
(IWS), University of Stuttgart, 70569 Stuttgart, Germany.}

\begin{article}
\footskip=12pt

\noindent\textbf{Contents of this file}
\begin{enumerate}
\item Texts S1--S5
\item Tables S1--S4
\end{enumerate}

\noindent\textbf{Introduction}

This Supporting Information provides the equations, constitutive assumptions,
boundary conditions, imaging details, comparison quantities, and numerical
checks that complement the main article. The organization follows the
Supporting Information of the related dynamic pore-network study, but the
equations below describe the continuum DuMu$^{\mathrm{X}}$ model used here; no
pore-throat invasion or pore-network upscaling relations are transferred.
Symbols follow the convention used in the main article: component or species
labels are superscripts, phase and location labels are subscripts, mass
fractions are denoted by uppercase $X$, and mole fractions by lowercase $x$.

Red author queries mark information that is not yet supported consistently by
the manuscript and the archived model provenance. They are not part of the
scientific claims and must be resolved before submission.

\sitext{Governing equations and primary variables}\label{app:governing}

The porous medium contains a wetting liquid phase $w$ and a gas phase $g$.
The mobile components are water ($w$), the gas carrier ($\mathrm{N_2}$ in the
experiment and the air component in the DuMu$^{\mathrm{X}}$ property package),
and dissolved NaCl. Local thermal equilibrium is assumed, so all phases and the
solid matrix share one temperature $T$. For component
$\kappa\in\{w,\mathrm{N_2},\mathrm{NaCl}\}$, the molar balance is
\begin{equation}
\frac{\partial}{\partial t}
\left[
  \phi\sum_{\alpha\in\{w,g\}}
  S_\alpha\rho_{\mathrm{mol},\alpha}x_\alpha^\kappa
  + c_{\mathrm{solid}}^\kappa
\right]
+\boldsymbol{\nabla}\!\cdot
\left[
  \sum_{\alpha\in\{w,g\}}
  \left(
  \rho_{\mathrm{mol},\alpha}x_\alpha^\kappa\boldsymbol{v}_\alpha
  +\boldsymbol{J}_\alpha^\kappa
  \right)
\right]
=q^\kappa,
\label{eq:si-component-balance}
\end{equation}
where $\phi$ is porosity, $S_\alpha$ is phase saturation,
$\rho_{\mathrm{mol},\alpha}$ is phase molar density,
$\boldsymbol{v}_\alpha$ is Darcy velocity,
$\boldsymbol{J}_\alpha^\kappa$ is the diffusive molar flux, and $q^\kappa$ is
an external source per bulk volume. The immobile storage term is zero for water
and the carrier gas. For NaCl it is
$c_{\mathrm{solid}}^{\mathrm{NaCl}}=
\phi_s\rho_{\mathrm{mol},s}^{\mathrm{NaCl}}$, where $\phi_s$ is the
precipitated-salt volume fraction per bulk volume. Precipitation therefore
transfers NaCl between dissolved and solid inventories without creating or
removing total NaCl.

Phase velocities are calculated with the multiphase extension of Darcy's law,
\begin{equation}
\boldsymbol{v}_\alpha
=-\frac{K k_{r,\alpha}}{\mu_\alpha}
\left(\boldsymbol{\nabla}p_\alpha-\rho_\alpha\boldsymbol{g}\right),
\label{eq:si-darcy}
\end{equation}
where $K$ is intrinsic permeability, $k_{r,\alpha}$ is relative permeability,
$\mu_\alpha$ is dynamic viscosity, $p_\alpha$ is phase pressure,
$\rho_\alpha$ is mass density, and $\boldsymbol{g}$ is gravitational
acceleration. Molecular diffusion is represented by the DuMu$^{\mathrm{X}}$
Fickian closure. In compact form,
\begin{equation}
\boldsymbol{J}_\alpha^\kappa
=-\phi S_\alpha\rho_{\mathrm{mol},\alpha}
D_{\alpha,\mathrm{eff}}^\kappa
\boldsymbol{\nabla}x_\alpha^\kappa,
\qquad
\sum_\kappa\boldsymbol{J}_\alpha^\kappa=\boldsymbol{0},
\label{eq:si-diffusion}
\end{equation}
where $D_{\alpha,\mathrm{eff}}^\kappa$ includes the porous-medium reduction of
the corresponding binary diffusion coefficient. The second relation enforces
the reference-component constraint; Equation~\ref{eq:si-diffusion} is therefore
not an additional independent flux law for every component.

The non-isothermal energy balance is
\begin{equation}
\frac{\partial u_{\mathrm{eff}}}{\partial t}
+\boldsymbol{\nabla}\!\cdot
\left[
\sum_{\alpha\in\{w,g\}}\rho_\alpha h_\alpha\boldsymbol{v}_\alpha
+\sum_{\alpha\in\{w,g\}}\sum_\kappa
h_\alpha^\kappa\boldsymbol{J}_\alpha^\kappa
-\lambda_{\mathrm{eff}}\boldsymbol{\nabla}T
\right]=q^{\mathrm{E}},
\label{eq:si-energy}
\end{equation}
where $u_{\mathrm{eff}}$ is the bulk internal-energy density,
$h_\alpha$ and $h_\alpha^\kappa$ are phase and component specific
enthalpies, $\lambda_{\mathrm{eff}}$ is effective thermal conductivity, and
$q^{\mathrm{E}}$ is an energy source. Advective enthalpy, diffusive enthalpy,
and heat conduction are retained. A separate crystallization-enthalpy source is
not imposed.

The finite-volume discretization uses the same conservative face flux for the
two materials at their vertical interface. Thus, interdomain exchange is an
outcome of the coupled pressure, composition, and temperature fields rather
than a prescribed transfer coefficient.

\sitext{Constitutive relations and salt equilibrium}\label{app:closures}

The saturations and capillary pressure satisfy
\begin{equation}
S_w+S_g=1,
\qquad
p_c=p_g-p_w.
\label{eq:si-phase-closure}
\end{equation}
For each porous material, the Brooks--Corey effective wetting saturation is
\begin{equation}
S_{w,e}=\frac{S_w-S_{w,r}}{1-S_{w,r}-S_{g,r}},
\label{eq:si-effective-saturation}
\end{equation}
and the implemented capillary pressure and relative permeabilities are
\begin{align}
p_c &= p_b S_{w,e}^{-1/\lambda_{\mathrm{BC}}},
\label{eq:si-bc-pc}\\
k_{r,w} &= S_{w,e}^{3+2/\lambda_{\mathrm{BC}}},
\label{eq:si-bc-krw}\\
k_{r,g} &= (1-S_{w,e})^2
\left(1-S_{w,e}^{1+2/\lambda_{\mathrm{BC}}}\right).
\label{eq:si-bc-krg}
\end{align}
Outside the physical interval, the numerical implementation regularizes or
clamps effective saturation as required by the nonlinear solver.

Liquid density, viscosity, enthalpy, and binary diffusion coefficients depend
on pressure, temperature, and composition through the DuMu$^{\mathrm{X}}$
brine--gas property package \cite{koch2021}. This salinity dependence is
essential because the density contrast is the modeled source of gravitational
instability. The liquid NaCl binary diffusivity retains the base property-model
value; no concentration-dependent hydrodynamic dispersion coefficient was
fitted in the simulations reported here.

The water activity is calculated from NaCl molality,
\begin{equation}
m_{\mathrm{NaCl}}
=\frac{x_w^{\mathrm{NaCl}}}{M^w x_w^w},
\label{eq:si-molality}
\end{equation}
where $M^w$ is the molar mass of water. The Pitzer-type osmotic coefficient
used in the archived implementation is \cite{Pitzer1973}
\begin{equation}
\phi_{\mathrm{osm}}(m)=
1-\frac{0.3915\sqrt{m}}{1+1.2\sqrt{m}}
+m\left(0.0765+0.2664\exp[-2\sqrt{m}]\right)
+0.00127m^2,
\label{eq:si-osmotic-coefficient}
\end{equation}
and
\begin{equation}
a_w=\exp\left[-2m_{\mathrm{NaCl}}M^w
\phi_{\mathrm{osm}}(m_{\mathrm{NaCl}})\right].
\label{eq:si-water-activity}
\end{equation}
The parameterization is evaluated at 293.15~K as a near-room-temperature
approximation; it is not fitted to the present evaporation experiment.

The solubility limit is the liquid NaCl mass fraction
$X_{w,\mathrm{max}}^{\mathrm{NaCl}}=0.26$. The corresponding saturation
molality is
\begin{equation}
m_{\mathrm{sat}}
=\frac{X_{w,\mathrm{max}}^{\mathrm{NaCl}}}
{M^{\mathrm{NaCl}}\left(1-X_{w,\mathrm{max}}^{\mathrm{NaCl}}\right)}
=6.012~\mathrm{mol\,kg^{-1}},
\label{eq:si-saturation-molality}
\end{equation}
where the value is rounded. Equilibrium precipitation and dissolution are
enforced through
\begin{equation}
c_{\mathrm{solid}}^{\mathrm{NaCl}}\geq0,
\qquad
g_{\mathrm{sat}}\geq0,
\qquad
c_{\mathrm{solid}}^{\mathrm{NaCl}}g_{\mathrm{sat}}=0,
\label{eq:si-complementarity}
\end{equation}
with
\begin{equation}
g_{\mathrm{sat}}=
\phi S_w\rho_{\mathrm{mol},w}
\left(m_{\mathrm{sat}}M^w x_w^w-x_w^{\mathrm{NaCl}}\right).
\label{eq:si-saturation-gap}
\end{equation}
This closure permits solid NaCl only when the dissolved phase is at saturation,
and permits undersaturation only when the solid inventory is zero.

Precipitated salt is stored as an immobile conserved inventory. In the
external-crust extension used in the manuscript, the inventory associated with
an atmospheric upper-boundary cell is interpreted as external salt exactly
once and is used to calculate an equivalent vapor resistance. It is not copied
to a second precipitation store. The underlying sand retains fixed $\phi$, $K$,
and retention functions, and no explicit crust-domain liquid flow is solved.
The numerical onset of a positive equilibrium solid inventory is therefore not
a prediction of nucleation time or crystal morphology.

\sitext{Initial and boundary conditions}\label{app:boundary}

\sisub{Upper evaporation boundary}

The water mass flux across the gas film is
\begin{equation}
j_{\mathrm{evap}}=
\frac{\rho_{g,f}D_g^{w,\mathrm{N_2}}}{\delta_g}
\left(X_{g,s}^{w,\mathrm{eq}}-X_{g,\mathrm{amb}}^w\right),
\label{eq:si-film-flux}
\end{equation}
where the gas properties are evaluated at representative film conditions and
$\delta_g$ is an effective film thickness. The equilibrium interfacial mole
fraction is
\begin{equation}
x_{g,s}^{w,\mathrm{eq}}
=\frac{a_w p_{\mathrm{sat}}(T_s)}{p_{\mathrm{atm}}},
\label{eq:si-vapor-equilibrium}
\end{equation}
and its corresponding mass fraction is
\begin{equation}
X_{g,s}^{w,\mathrm{eq}}=
\frac{M^w x_{g,s}^{w,\mathrm{eq}}}
{M^w x_{g,s}^{w,\mathrm{eq}}+
M^{\mathrm{N_2}}\left(1-x_{g,s}^{w,\mathrm{eq}}\right)}.
\label{eq:si-mole-to-mass}
\end{equation}
Dry nitrogen is represented by $X_{g,\mathrm{amb}}^w=0$. The signed form of
Equation~\ref{eq:si-film-flux} permits condensation in principle, although it
does not occur under the prescribed dry-gas condition.

For an upper boundary face, conserved external NaCl mass $M_{\mathrm{cr}}$ is
converted to an equivalent porous-crust thickness,
\begin{equation}
h_{\mathrm{cr}}=
\frac{M_{\mathrm{cr}}}
{\rho_{\mathrm{cr}}(1-\varepsilon_{\mathrm{cr}})A},
\label{eq:si-crust-thickness}
\end{equation}
where $\rho_{\mathrm{cr}}$ is crystal density,
$\varepsilon_{\mathrm{cr}}$ is assumed crust porosity, and $A$ is the native
boundary-face area. Series diffusion through the gas film and equivalent crust
then gives
\begin{equation}
j_{\mathrm{evap,cr}}=
\frac{j_{\mathrm{evap}}}
{1+\displaystyle\frac{h_{\mathrm{cr}}D_g^{w,\mathrm{N_2}}}
{\delta_gD_{\mathrm{cr}}^w}}.
\label{eq:si-crust-flux}
\end{equation}
The exploratory coupled case uses $\rho_{\mathrm{cr}}=2165~\mathrm{kg\,m^{-3}}$,
$\varepsilon_{\mathrm{cr}}=0.43$, complete within-face coverage, and
$D_{\mathrm{cr}}^w/D_g^{w,\mathrm{N_2}}=0.08$. These values are conditional
assumptions, not independently measured properties of the observed moist
efflorescence. Additional crust thermal resistance and explicit hydraulic
transport inside the external salt layer are excluded.

The upper energy flux combines water-vapor enthalpy, carrier-gas enthalpy, and
gas-film conduction:
\begin{equation}
j_{\mathrm{E},s}=
h_g^w j_{\mathrm{evap}}
+h_g^{\mathrm{N_2}}j_g^{\mathrm{N_2}}
+\frac{\lambda_g}{\delta_g}(T_s-T_{\mathrm{amb}}).
\label{eq:si-top-energy}
\end{equation}
The same effective thickness is used for mass and heat transfer.

\sisub{Reservoir connection and external walls}

The reservoir pressure at boundary elevation $z_b$ is
\begin{equation}
p_{w,\mathrm{res}}(z_b)=p_{\mathrm{atm}}
+\rho_{w,\mathrm{res}}g(z_{\mathrm{res}}-z_b),
\label{eq:si-reservoir-pressure}
\end{equation}
where $z_{\mathrm{res}}$ is 20~mm above the sand--filter interface. The central
inlet exchanges liquid with the explicitly resolved filter according to the
hydraulic-potential difference between the reservoir and the adjacent filter
cell. Denoting boundary transmissibility by $\mathcal{T}_b$,
\begin{equation}
q_{V,b}=-\mathcal{T}_b
\left(\Psi_{w,b}-\Psi_{w,\mathrm{res}}\right).
\label{eq:si-bottom-flux}
\end{equation}
The sign determines the upstream state used for component fractions and
enthalpy; both inflow and return flow are admitted. No second Robin-type filter
resistance is applied in addition to the explicit filter.

All other external boundaries are impermeable to fluid and components. Heat
exchange through the cylindrical and bottom holder walls is
\begin{equation}
j_{\mathrm{E},\mathrm{wall}}=
\frac{\lambda_{\mathrm{wall}}}{d_{\mathrm{wall}}}
(T-T_{\mathrm{amb}}),
\label{eq:si-wall-heat}
\end{equation}
using $\lambda_{\mathrm{wall}}=0.184~\mathrm{W\,m^{-1}\,K^{-1}}$ and
$d_{\mathrm{wall}}=5~\mathrm{mm}$.

\sisub{Initial state and pre-equilibration}

The domain is initialized fully saturated with 1~mol~L$^{-1}$ NaCl solution at
293.15~K. Before evaporation, the numerical column is pre-equilibrated for 2~h
with the gas-film water flux inactive. Reported simulation time $t=0$
corresponds to the start of evaporation. Geometry and boundary parameters are
summarized in Table~\ref{tab:si-domain}; material parameters are listed in
Table~\ref{tab:si-materials}.

\verify{The main article currently states that $\delta_g$ was calibrated using
only the first 24~h, whereas archived provenance for the selected
$\delta_g=8.100162$~mm value documents a full 174-h evaporation calibration.
The calibration window and the status of 24--174~h as independent validation
must be reconciled before submission.}

\sitext{MRI acquisition, calibration, processing, and uncertainty}
\label{app:mri}

The \NaIso{} and \HIso{} acquisition parameters are reproduced in
Table~\ref{tab:si-mri}. After the initial \NaIso{} scan, the schedule was
$8\times30$~min, $2\times60$~min, $3\times120$~min, and
$27\times360$~min, for a total experiment duration of 174~h. Each image is
assigned its nominal acquisition time. The finite 32-min \NaIso{} acquisition
window is not a point measurement and should be considered when interpreting
rapid early-time redistribution.

For calibration, homogeneous F36 and W3 packings were saturated with 0, 1, 2,
3, 4, and 5~mol~L$^{-1}$ solutions and scanned using the same protocol. The
Na volume concentration is
\begin{equation}
c_{\mathrm{vol}}^{\mathrm{Na}}
=\frac{n^{\mathrm{Na}}}{V_{\mathrm{bulk}}}
=c_{\mathrm{sol}}^{\mathrm{Na}}\theta_{\mathrm{vol}},
\label{eq:si-mri-calibration}
\end{equation}
where $n^{\mathrm{Na}}$ is the amount of Na, $V_{\mathrm{bulk}}$ is bulk
volume, $c_{\mathrm{sol}}^{\mathrm{Na}}$ is solution concentration, and
$\theta_{\mathrm{vol}}$ is volumetric water content. A linear response was used
for F36 and a second-order response for W3:
\begin{align}
c_{\mathrm{vol,F36}}^{\mathrm{Na}} &=
a_{\mathrm{F36}}S_r+b_{\mathrm{F36}},
\label{eq:si-f36-calibration}\\
c_{\mathrm{vol,W3}}^{\mathrm{Na}} &=
a_{\mathrm{W3}}S_r^2+b_{\mathrm{W3}}S_r+d_{\mathrm{W3}},
\label{eq:si-w3-calibration}
\end{align}
where $S_r$ is the processed relative voxel signal.

\verify{Insert the fitted coefficients, their units, fit residuals, and the
number of independent calibration samples for
Equations~\ref{eq:si-f36-calibration} and \ref{eq:si-w3-calibration}. A plot
alone is insufficient for reproducibility.}

Raw k-space data were reconstructed by inverse Fourier transformation in
Paravision~6. Subsequent processing in ImageJ~1.53t used common display scaling,
a 1.5-pixel median filter, and the volume-viewer plugin for transparent 3-D
renderings. A spike artifact within the sample region was suppressed for the
3-D visualization by filtering high voxel values. Quantitative profiles must be
distinguished from display-only brightness scaling and artifact suppression.

The standard deviations obtained by back-calculating the six calibration
solutions were 0.11, 0.26, 0.35, 0.42, 0.51, and
0.63~mol~L$^{-1}$ for F36, and 0.13, 0.28, 0.29, 0.54, 0.67, and
0.63~mol~L$^{-1}$ for W3. These correspond to mean relative errors of about
14\% for F36 and about 15\% for W3 at the higher concentrations; relative error
is larger near zero concentration. These voxel variations quantify repeatability
within the calibration images under the adopted processing, not a complete
voxelwise confidence interval. They exclude fitted-coefficient uncertainty,
finite acquisition time, spatial registration, interface partial volumes,
undetected precipitated NaCl, and possible supersaturation beyond the
calibration range.

The \HIso{} data were normalized to the initial scan using an initial
volumetric water content of 0.38 because a separate marker tube could not be
accommodated. The \HIso{} series is therefore best interpreted as a relative
check on water-content changes, not as an independently calibrated absolute
water-content measurement at every voxel.

\sitext{Comparison quantities and numerical verification status}
\label{app:verification}

For subdomain $\Omega_i$ (F36, W3, or the complete sample), the amount of
MRI-visible dissolved sodium is
\begin{equation}
n_{\Omega_i}^{\mathrm{Na}}(t)=
\int_{\Omega_i}c_{\mathrm{vol}}^{\mathrm{Na}}(\boldsymbol{x},t)
\,\mathrm{d}V.
\label{eq:si-integrated-na}
\end{equation}
The simulated dissolved NaCl field is converted to the corresponding Na amount
before comparison. Simulated solid NaCl is excluded because \NaIso-MRI does not
detect the solid inventory.

Let $\Gamma_{\mathrm{FW}}$ denote the vertical interface with unit normal from
F36 to W3. The conservative NaCl transfer rate is
\begin{equation}
\dot n_{\mathrm{FW}}^{\mathrm{NaCl}}=
\int_{\Gamma_{\mathrm{FW}}}
\left(
\rho_{\mathrm{mol},w}x_w^{\mathrm{NaCl}}\boldsymbol{v}_w
+\boldsymbol{J}_w^{\mathrm{NaCl}}
\right)\!\cdot\boldsymbol{n}_{\mathrm{FW}}\,\mathrm{d}A,
\label{eq:si-interface-flux}
\end{equation}
with the two terms reported separately as advective and diffusive contributions.
Positive values indicate net transfer from F36 to W3.

For $N$ matched profile values, concentration error is
\begin{equation}
\mathrm{RMSE}_c=
\left[
\frac{1}{N}\sum_{i=1}^{N}
\left(c_{\mathrm{sim},i}^{\mathrm{Na}}-
c_{\mathrm{MRI},i}^{\mathrm{Na}}\right)^2
\right]^{1/2}.
\label{eq:si-rmse}
\end{equation}
RMSE is reported separately for F36 and W3 because pooling both halves can mask
opposite biases. The sampling volume, depth interval, spatial averaging, and
nominal acquisition time must accompany every value.

Discrete conservation is checked from accepted time steps. For component
$\kappa$, the cumulative residual is
\begin{equation}
r^\kappa(t)=N^\kappa(t)-N^\kappa(0)
+\int_0^t\dot N_{\partial\Omega}^\kappa(\tau)\,\mathrm{d}\tau,
\label{eq:si-budget-residual}
\end{equation}
where outward boundary flow is positive. The selected 8.100-mm film case
completed 2~h pre-equilibration plus 174~h evaporation. Its maximum relative
water, salt, and material budget residuals were of order $10^{-12}$; the maximum
relative energy residual was $1.64\times10^{-10}$. These values establish
internal numerical accounting, not physical validity.

The fixed-geometry case without additional crust resistance predicted reservoir
solution loss of 52.4057~g compared with 51.52~g measured, an endpoint
difference of +1.72\%. This did not meet the pre-registered 1\% endpoint
tolerance. The exploratory coupled external-crust case predicted 51.5403~g
(+0.0395\% at 174~h), but this close terminal match does not independently
validate the crust mechanism because the crust properties and mapping are
assumed and the gas-film thickness was already calibrated. Short hydrostatic, boundary-transport,
conservation, sparse-output, and thread-equivalence gates passed. However, a
targeted radial refinement of the approximately 4-mm inlet and a complete
production-scale grid/time-step convergence study remain unresolved. The inlet
spans only about 1.6 maximum bottom-face equivalent diameters in the current
mesh. Only saturated filter permeability is measured; filter porosity,
retention, relative permeability, and thermal properties remain provisional.
Table~\ref{tab:si-verification} separates passed checks from open limitations.

\verify{Do not describe the model as \emph{validated} until the gas-film
calibration window is reconciled and an independent prediction interval is
defined. Numerical closure, thread equivalence, and completion to 174~h do not
substitute for grid convergence or independent experimental validation.}

\begingroup
\jdraftfalse
\normalsize
\bibliography{literature}
\endgroup

\end{article}
\clearpage

\begin{table}
\settablenum{S1}
\caption{Principal geometry, initial conditions, and boundary parameters.}
\label{tab:si-domain}
\centering
\begin{tabular}{p{0.47\textwidth}p{0.23\textwidth}p{0.22\textwidth}}
\toprule
Quantity & Value & Status \\
\midrule
Sand-column inner diameter & 50~mm & experimental geometry \\
Sand height & 40~mm & experimental geometry \\
Porous-filter thickness & 4.5~mm & experimental geometry \\
Filter saturated permeability & $5.30\times10^{-12}$~m$^2$ & measured mean \\
Central inlet diameter & approximately 4~mm & approximate geometry \\
Reservoir level above sand--filter interface & 20~mm & experimental setting \\
Initial/reservoir NaCl concentration & 1~mol~L$^{-1}$ & experimental setting \\
Initial, reservoir, and ambient temperature & 293.15~K & prescribed \\
Atmospheric pressure & $1.00\times10^5$~Pa & prescribed \\
Ambient relative humidity & 0 & dry-gas assumption \\
Effective gas-film thickness $\delta_g$ & 8.100162~mm & calibration window to reconcile \\
Numerical pre-equilibration & 2~h & evaporation inactive \\
Maximum time-step size & 300~s & numerical control \\
Wall conductivity $\lambda_{\mathrm{wall}}$ & 0.184~W~m$^{-1}$~K$^{-1}$ & prescribed \\
Wall thickness $d_{\mathrm{wall}}$ & 5~mm & prescribed \\
NaCl solubility mass fraction $X_{w,\mathrm{max}}^{\mathrm{NaCl}}$ & 0.26 & equilibrium closure \\
Crust crystal density $\rho_{\mathrm{cr}}$ & 2165~kg~m$^{-3}$ & assumed \\
Equivalent-crust porosity $\varepsilon_{\mathrm{cr}}$ & 0.43 & assumed \\
Crust/gas diffusivity ratio & 0.08 & sensitivity choice \\
Within-face crust coverage & 1 & assumed complete coverage \\
\bottomrule
\end{tabular}
\end{table}

\begin{table}
\settablenum{S2}
\caption{Hydraulic material parameters used for F36 and W3. The distinction
between reported saturated water content and computational porosity is retained.}
\label{tab:si-materials}
\centering
\begin{tabular}{p{0.47\textwidth}cc}
\toprule
Quantity & F36 & W3 \\
\midrule
Mean particle size ($\mu$m) & 160 & 90 \\
Particle-size range ($\mu$m) & 90--355 & 2--400 \\
Intrinsic permeability $K$ (m$^2$) & $7.83\times10^{-12}$ & $3.20\times10^{-14}$ \\
Reported saturated water content $\theta_{\mathrm{sat}}$ & 0.39 & 0.31 \\
Reported residual water content $\theta_{\mathrm{res}}$ & 0.02 & 0.08 \\
Computational porosity $\phi$ & 0.40 & 0.36 \\
Computational residual saturation $S_{w,r}$ & 0.05 & 0.22 \\
Residual gas saturation $S_{g,r}$ & 0 & 0 \\
Entry pressure $p_b$ (Pa) & $4.69\times10^3$ & $1.12\times10^4$ \\
Brooks--Corey exponent $\lambda_{\mathrm{BC}}$ & 5.04 & 0.384 \\
\bottomrule
\end{tabular}

\vspace{1ex}
\noindent\textcolor{red}{\textbf{Author query---}The main article reports
$\theta_{\mathrm{sat}}=0.39/0.31$, while current computational inputs use
$\phi=0.40/0.36$. Confirm the provenance and interpretation of both pairs,
especially W3, rather than treating them as interchangeable.}
\end{table}

\begin{table}
\settablenum{S3}
\caption{MRI acquisition parameters used during the heterogeneous experiment.}
\label{tab:si-mri}
\centering
\begin{tabular}{lcc}
\toprule
Parameter & \NaIso{} imaging & \HIso{} imaging \\
\midrule
Pulse sequence & 3DSE & MSME \\
$t_R$ (ms) & 300 & 3000 \\
$t_E$ (ms) & 4.6 & 4 \\
Slices & 1 & 40 \\
Slice thickness (mm) & 55 & 1 \\
Number of averages & 4 & 4 \\
Scan time (min) & 32 & 51 \\
Matrix size & $64\times40\times40$ & $256\times256\times1$ \\
Resolution (mm) & $1\times1.375\times1.375$ & $0.25\times0.25\times0.25$ \\
Bandwidth (kHz) & 50 & 150 \\
\bottomrule
\end{tabular}
\end{table}

\begin{table}
\settablenum{S4}
\caption{Status of numerical checks and remaining evidence limitations for the
selected 8.100-mm gas-film baseline and the exploratory external-crust extension.}
\label{tab:si-verification}
\centering
\small
\begin{tabular}{>{\raggedright\arraybackslash}p{0.28\textwidth}
                >{\raggedright\arraybackslash}p{0.25\textwidth}
                >{\raggedright\arraybackslash}p{0.39\textwidth}}
\toprule
Check & Status & Interpretation \\
\midrule
174-h completion & passed & Solver completion only; not model validation \\
Water/salt/material budgets & residuals $\sim10^{-12}$ & Conservative accounting on accepted steps \\
Energy budget & residual $1.64\times10^{-10}$ & Internal energy accounting \\
Hydrostatic and bidirectional boundary tests & passed & Implementation consistency \\
Thread equivalence (1/2/4/8 threads) & passed for registered benchmark & Reproducibility of tested configuration \\
Reservoir-loss endpoint & +1.72\% error & Did not meet pre-registered 1\% tolerance \\
External-crust endpoint & +0.0395\% error & Close fit is conditional, not mechanism validation \\
Gas-film identification & unresolved provenance & First-24-h and full-window claims conflict \\
Local inlet radial refinement & open & Current aperture is locally under-resolved \\
Production grid/time-step convergence & open & No general convergence claim is supported \\
Filter unsaturated/thermal properties & provisional & Saturated permeability alone is measured \\
Sand geometry feedback & excluded & Porosity and permeability remain fixed \\
Explicit crust hydraulics/thermal resistance & excluded & Only vapor resistance is represented \\
\bottomrule
\end{tabular}
\end{table}

\end{document}





